\section{Conclusion}\label{sec:conclusion}

The proof in this paper has a shape that we expect to survive extensions of the calculus. A new sort of member, whether a path, a capture set, or a classifier, adds constructors to types and propositions, evidence rules for them, and cases to the translation. The normalizer gains cases, the canonical forms theorem gains cases, and preservation and progress do not change, because they only ever ask the theorem for a shape. Nothing in the argument depends on inverting a derivation of the source calculus. We have not carried out such an extension, and the claim is an argument from the structure of the proof, not a theorem.

For paths the changes are known. Paths gain field selection and singleton types, blocks become keyed by paths rather than variables, and the translation let-binds every path prefix that a type mentions, so that a target block exists for it. Singleton evidence between a variable and a path is then an equality between blocks. This is the one place where the translation stops being homomorphic on types.

Two further uses of the calculus follow from its being an explicit-evidence intermediate language. A small surface language with a derivation-producing type checker would let programs be written directly, with the \FCdot checker validating the evidence the front end produced. And a compiler for a language in the DOT family could emit evidence for the subsumption steps it takes, in the way GHC emits FC coercions, and have them validated independently of its own subtyping algorithm. The second is where we expect the design to matter most, for the parts of a compiler that have no reference semantics beyond a paper.
